% =====================================================
% 题型：解三角形多选题 (q_tjw_20260605_07_triangle_identities_3.tex)
% =====================================================
% =====================================================
% 难度：★★ (中档)
% =====================================================
\newcommand{\qTriangleIdentitiesThreeStem}{%
  在 $\triangle ABC$ 中，内角 $A, B, C$ 所对的边分别为 $a, b, c$，如下判断正确的是%
}

\newcommand{\qTriangleIdentitiesThreeOptions}{%
  \mcoptions[1]{
    若 $A > B$，则 $\cos A > \cos B$
  }{
    若 $b\cos A - a\cos B = 0$，则 $\triangle ABC$ 为等腰三角形
  }{
    若 $\cos A\cos B\cos C > 0$，则 $\triangle ABC$ 为锐角三角形
  }{
    若满足条件 $A = \dfrac{\pi}{4}$，$c = 2$ 的 $\triangle ABC$ 有两个，则 $a$ 的取值范围为 $(\sqrt{2}, 2)$
  }%
}

\newcommand{\qTriangleIdentitiesThreeAnswer}{B, C, D}

\newcommand{\qTriangleIdentitiesThreeSolution}{%
  对各选项逐一进行判定：
  
  对于 \textbf{A} 选项：
  因为 $y = \cos x$ 在 $(0, \pi)$ 上是单调递减的。
  若 $A > B$，则 $\cos A < \cos B$，故 \textbf{A} 错误；
  
  对于 \textbf{B} 选项：
  由正弦定理边化角得：
  \[ \sin B\cos A - \sin A\cos B = 0 \implies \sin(B-A) = 0. \]
  因为 $A, B$ 是三角形内角，且 $B-A \in (-\pi, \pi)$，所以 $B-A = 0 \implies A = B$。
  所以 $\triangle ABC$ 是等腰三角形，故 \textbf{B} 正确；
  
  对于 \textbf{C} 选项：
  因为在任意三角形中，最多只能有一个钝角（即最多只有一个角的余弦值为负）。
  若 $\cos A\cos B\cos C > 0$，说明三个余弦值全部为正数，即三个内角全部为锐角。
  所以 $\triangle ABC$ 为锐角三角形，故 \textbf{C} 正确；
  
  对于 \textbf{D} 选项：
  已知 $A = \dfrac{\pi}{4}$，$c = 2$，若三角形有两个解，则需满足 $c\sin A < a < c$。
  代入已知数据有：
  \[ 2\sin 45^{\circ} < a < 2 \implies \sqrt{2} < a < 2. \]
  所以 $a$ 的取值范围是 $(\sqrt{2}, 2)$，故 \textbf{D} 正确。
  
  综上所述，正确选项为 \textbf{B, C, D}。
}

\newcommand{\qTriangleIdentitiesThreeDiagnosis}{%
  用于课堂精准变式训练。考查余弦函数的单调性、两角差的正弦公式、三角形内角的钝角唯一性、以及三角形解的个数边界值确定。
}

\newcommand{\qTriangleIdentitiesThreeTeachingNotes}{%
  \item \textbf{单调性排雷}：A 选项是极其容易因马虎而做错的。大边对大角，但余弦函数是单调递减的，所以在三角形中大角对应的余弦值更小。
  \item \textbf{乘积符号判定}：C 选项帮助学生理解余弦乘积符号在确定三角形形状时的应用。
}


% =====================================================
% 别名兼容：旧课次宏定义
% =====================================================
\newcommand{\qTJWTriangleIdentitiesThreeStem}{\qTriangleIdentitiesThreeStem}
\newcommand{\qTJWTriangleIdentitiesThreeOptions}{\qTriangleIdentitiesThreeOptions}
\newcommand{\qTJWTriangleIdentitiesThreeAnswer}{\qTriangleIdentitiesThreeAnswer}
\newcommand{\qTJWTriangleIdentitiesThreeSolution}{\qTriangleIdentitiesThreeSolution}
\newcommand{\qTJWTriangleIdentitiesThreeDiagnosis}{\qTriangleIdentitiesThreeDiagnosis}
\newcommand{\qTJWTriangleIdentitiesThreeTeachingNotes}{\qTriangleIdentitiesThreeTeachingNotes}
